% !TEX TS-program = lualatex+se
% !TEX encoding = UTF-8

\documentclass[oneside, 12pt, a4paper]{book}

% --- USTAWIENIA STRONY ---
\usepackage{geometry}
\geometry{
%	showframe=true,
%	showcrop=true,
	left=25mm,
	right=20mm,
	top=20mm,
	bottom=18mm,
	marginparwidth=0mm,
	marginparsep=0mm
	}
%\usepackage[letterpaper,showframe]{geometry}	
% usuwanie nagłówków/stopek z\,dokumentu
\usepackage[indent=0pt]{parskip} % ustawianie paragrafów
%\setstretch{1.5} % półtorej linii odległości między wierszami
\usepackage{titlesec} % ustawianie nagłówków
%\usepackage{fullpage} % redukcja marginesów
\usepackage{setspace} % odstępy między liniami
\usepackage{enumitem}% listy numerowane
% -----------------------------

%----------- style nagłówków ----------------
\newcommand{\headingone}{\scshape\Large}
\newcommand{\headingtwo}{\large}
\newcommand{\margintext}[1]{\textit{\small #1}}
% przesunięcia poziome
\newcommand{\smallwidth}{\hskip 1em}
\newcommand{\medwidth}{\hskip 3pt}
\newcommand{\largewidth}{\hskip 12pt}
% przesunięcia pionowe
\newcommand{\verspace}{0.3}
\newcommand{\linespace}{\vskip 6pt}
\newcommand{\paragraphspace}{\vskip 10pt}
\newcommand{\segmentspace}{\vskip 16pt}
\newcommand{\sectionspace}{\vskip 28pt}
%------------------------------------------

% --- GREGORIO ---	
\usepackage{gregoriotex}
\gresetlinecolor{gregoriocolor}
\gresetheadercapture{annotation}{greannotation}{string}% odczytuje adnotacje nad inicjałem jako ciąg, co pozwala na jego renderowanie
\gresetheadercapture{commentary}{grecommentary}{string}% odczytuje komentarz jako ciąg, co pozwala na jego renderowanie
\grechangedim{commentaryraise}{0.3cm}{scalable} % podnosi komentarz wystarczająco ponad górną linię
\grechangestyle{annotation}{\color{black}\small}% format adnotacji
\grechangestyle{commentary}{\color{gregoriocolor}\small}% format komentarza
\grechangestyle{abovelinestext}{\color{gregoriocolor}\small}% format napisu nad liniami
\grechangestyle{initial}{\fontsize{36}{36}\selectfont} % zmniejsza wielkość inicjału do bardziej proporcjonalnego rozmiaru (=36pt)
\grechangedim{intersyllablespacenotes}{0.26cm}{scalable}% Minimalna odległość między nutami różnych sylab.

% umieszczenie zawartości składu ich unikalnego argumentu po lewej i prawej stronie każdej linii
\let\grelocalleftbox\localleftbox
\let\grelocalrightbox\localrightbox

%\newcommand{\separator}{\greseparator{3}{40}}
%\newfontface\vectis{VectisMonumentalW01-Regular}
%\def\greinitialformat#1{\fontsize{20}{20}\vectis #1}
% ------------------

%\usepackage[hyperindex=true]{hyperref}

%---- usun konflikt między paracol i luacolor ----
%-------- kod poprawiający konflikt między paracol i luacolor ----
\ifluatex
\usepackage{luacolor}
\directlua{
 resetboxcolor = function (head,color)
  for n in node.traverse(head) do
    if 
      node.has_attribute(n,oberdiek.luacolor.getattribute(),
        oberdiek.luacolor.getvalue(token.get_macro("default@color")))
    then 
      node.set_attribute(n,oberdiek.luacolor.getattribute(),color)      
    end
    if n.id == node.id("hlist") or n.id == node.id("vlist") then
     resetboxcolor (n.list,color)
    end
 end
 return head
 end 
}

\protected\def\luacolorresetboxcolor#1{%
  \directlua {local color= tex.getattribute(oberdiek.luacolor.getattribute())
              resetboxcolor (tex.getbox(\the#1),color)}}

\else
\protected\def\luacolorresetboxcolor#1{}
\fi
%------ koniec kodu --------------------------------------
%-----------------------------------------------

% --- SPIS TREŚCI -------
%\usepackage{imakeidx}
%\makeindex[title=Alfabetyczny spis program=texindy,options=-M styl.xdy]
%\usepackage{filecontents}
%\begin{filecontents*}{styl.xdy}
%;;; xindy style file
%(markup-locclass-list :open "\dotfill" :sep "")
%(define-attributes (( "textbf" "default" )) )
%(markup-locref   :attr  "textbf"     :open "\textbf{\hyperpage{" :close "}}")
%(markup-locref   :attr  "textit"     :open "\textit{\hyperpage{" :close "}}")
%(markup-locref   :attr  "textttt"     :open "\textttt{\hyperpage{" :close "}}")
%(markup-locref   :attr  "texttsc"     :open "\texttsc{\hyperpage{" :close "}}")
%(markup-locref   :attr  "default"     :open "\hyperpage{" :close "}")
%\end{filecontents*}
% -------------------------

% ---  Konfiguracja kolumn ---
\usepackage{paracol}
\setcolumnwidth{.49\textwidth, .49\textwidth}
\setlength{\columnseprule}{0.4pt}
\setlength{\columnsep}{.02\textwidth}
\colseprulecolor{red}% wymagany pakiet xcolor
\footnotelayout{m}
% ----------------------------

\makeatletter 
\patchcmd\pcol@hfil
 {\copy\pcol@tempboxa}
 {\luacolorresetboxcolor\pcol@tempboxa\copy\pcol@tempboxa}{}{\fail}
\makeatother
% ----------------------------

\footnotelayout{m}

%% --- UKŁAD DWUKOLUMNOWY Z TŁUMACZENIEM ---
%\usepackage{tcolorbox}
%\newtcolorbox{nutybox}{size=minimal, colback=white, colframe=red, arc=0mm, width=\linewidth, halign=left, valign=center}
%\newtcolorbox{tlumbox}{size=minimal, colback=white, colframe=red, arc=0mm, width=\linewidth, halign=left, valign=top, fontupper=\small\itshape} 
%\newcommand{\nuty}[1]{\begin{nutybox}#1 \end{nutybox}}
%\newcommand{\tlum}[1]{\begin{tlumbox}\tlumacz{#1} \end{tlumbox}}
%% ------------------------------------------------------

\usepackage{Tabbing}
\usepackage{longtable}
\usepackage[savepos]{zref}

% --- CZCIONKI DOKUMENTU ---
\usepackage[latin,polish]{babel}
\usepackage{fontspec}
%\usepackage{times}
%\usepackage[math]{anttor}
%\usepackage{tgtermes}
%\usepackage{mathptmx}
\usepackage{ebgaramond}
\usepackage{libertine}
\defaultfontfeatures{Ligatures=TeX}

%\setmainfont{TeX Gyre Termes}
%\setmainfont{TeX Gyre Heros}
%\setmainfont{Calluna Regular}
\setmainfont{EB Garamond}
%\setmainfont{Bodoni Moda}
%\newfontfamily{\sekfont}{Times New Roman}
\newfontfamily{\sekfont}{Libertinus Serif}
\newcommand\sekcjafont{\fontsize{16}{16}\selectfont \sekfont}

%----- Strona tytułowa -------------
\newfontfamily{\tytfont}{JSL Ancient}
\newcommand\tytulfont[1]{\fontsize{36}{30}\selectfont \tytfont#1\par}
\newcommand\podtytulfont[1]{\fontsize{18}{10}\selectfont \tytfont#1\par}

%\setmainfont[<font features>]{<name of font>} % sets the roman font
%\setsansfont[<font features>]{<name of font>} % sets the sans font
%\setmonofont[<font features>]{<name of font>} % sets the monospace font

% --- KONIEC CZCIONKI DOKUMENTU ---

% --- GRAFIKA I KOLORY ---
\usepackage{xcolor} % kolory
\usepackage{graphicx} % umożliwia import obrazów
\graphicspath{{img/}} % katalog obrazów
\usepackage{pgfornament}% ornamenty graficzne
% --- KONIEC GRAFIKA I KOLORY -----------

% --- KOMENDY WŁASNE ---
\newcommand{\red}[1]{\textcolor{red}{#1}}
\newcommand{\black}[1]{\textcolor{black}{#1}}
% --------------------------------
\newcommand{\tlumacz}[1]{\textit{#1}}
\newcommand{\inicjal}[2]{\textcolor{red}{\LARGE\textbf{#1}}{#2}}
%--------------------------------
\titlespacing*{\chapter}{0pt}{0pt}{5pt}	
\titleformat{\chapter}[hang]{\filcenter\LARGE\scshape\bfseries}{}{0pt}{} % format nagłówków
\titlespacing*{\section}{0pt}{0pt}{0pt}	
\titleformat{\section}[hang]{\filcenter\large\scshape}{}{0pt}{} % format nagłówków
% ----------- RUBRYKA -------------------
\renewcommand\fbox{\fcolorbox{red}{white}}

\newcommand{\rubryka}[1]{
	\begin{spacing}{0.8}
	\begin{center}
		\parbox{0.9\textwidth}{\centering\textcolor{red}{\small\textit{#1}}}
	\end{center}
	\end{spacing}
	\linespace}
% ----------- SEKCJA --------------------
\newcommand{\sekcja}[2]{
	\begin{spacing}{0.6}
	\begin{center}
		\section*{#1}
		\linespace
		\begin{minipage}{0.88\linewidth}
			\red{\small\textit{#2}}
		\end{minipage}
	\end{center}	
	\end{spacing}
	\paragraphspace}
%------------DEKOR ---------------------
\newcommand{\dekor}[2]{
	\begin{center}
		\vspace*{\fill}
		\includegraphics[width=#1\textwidth]{#2}
		\vspace*{\fill}
	\end{center}}
%------- Zarządzanie wersami w psalmach z tlumaczeniem -------------------------
\newcounter{licznik}
\newcommand{\licznikZero}{\setcounter{licznik}{0}}% zeruj licznik
\newcommand{\licznikUp}{\refstepcounter{licznik}}% zwiększ licznik 
\licznikZero
\newcommand{\BeginParallel}{\begin{spacing}{0.9} \begin{paracol}{2} \licznikZero}
\newcommand{\lat}[1]{\selectlanguage{latin}\ensurevspace{2\baselineskip}\begin{minipage}[t]{\linewidth}\ifnum \value{licznik}>0 \licznikUp \thelicznik .~{#1}\par \else \licznikUp{#1}\par \fi \end{minipage}\switchcolumn}
\newcommand{\latN}[1]{\selectlanguage{latin}\ensurevspace{2\baselineskip}\begin{minipage}[t]{\linewidth}{#1}\par\end{minipage}\switchcolumn}% kolumna bez numeracji wersów
\newcommand{\pol}[1]{\selectlanguage{polish}\begin{minipage}[t]{\linewidth}{#1}\par\end{minipage}\switchcolumn*}
\newcommand{\EndParallel}{\end{paracol} \end{spacing}}

% --- ZNAKI SPECJALNE ---
% 
%\newcommand †
% --------------------------
\begin{document}
%\begin{center}
%	\includegraphics[width=1.0\textwidth]{Purification_BVM02.jpg}
%	\vfill
%	\podtytulfont{AD I. VESPERAS}
%	\tytulfont{In Purificatione\\Beatæ Mariæ Virginis}
%	\vfill
%	\large\textit{Liber usualis, Solesmes, 1961, p.1355}
%\end{center}
%\newpage
%% ----------------- Modlitwy wstępne -----------
\sekcja{Tłumaczenia\footnotemark}{}
%---------------------
%\small
\vspace{-7mm}
\red{\textbf{Ant.1}} O cudowna zamiano: * Stworzyciel rodzaju ludzkiego przyjmując ciało z duszą, raczył się narodzić z Dziewicy: i stając się człowiekiem bez udziału człowieka, udzielił nam swego Bóstwa.

\red{\textbf{Psalm 109}}\\
1. Rzekł Pan Panu memu: * Siądź po prawicy mojej:\\
2. Aż położę nieprzyjaciół twoich * podnóżkiem nóg twoich.\\
3. Berło mocy twojej wypuści Pan z Syjonu: * panuj wpośród nieprzyjaciół twoich.\\
4. Przy tobie panowanie w dzień mocy twej w blaskach świętości: * z żywota przed jutrzenką zrodziłem cię.\\
5. Przysiągł Pan i nic będzie żałował: * Ty jesteś kapłanem na wieki, według porządku Melchizedecha.\\
6. Pan po prawicy twojej, * poraził w dzień gniewu swego królów.\\
7. Będzie sądził narody, dopełni zagłady, * potłucze głowy wielu na ziemi.\\
8. Ze strumienia na drodze pić będzie: * dlatego wywyższy głowę.

%---------------------
\red{\textbf{Ant.2}} Kiedyś się narodził * w niewymowny sposób z Dziewicy, wypełniły się Pisma. Zstąpiłeś jako rosa na runo, aby zbawić rodzaj ludzki! Ciebie chwalimy, Boże nasz.

\red{\textbf{Psalm 112}}\\
1. Chwalcie, słudzy, Pana, * chwalcie imię Pańskie.\\
2. Niech będzie imię Pańskie błogosławione * odtąd i aż na wieki.\\
3. Od wschodu słońca aż do zachodu, * chwalebne jest imię Pańskie.\\
4. Wywyższony jest Pan nad wszystkie narody * i nad niebiosa chwała Jego.\\
5. Któż jest jak Pan, Bóg nasz, który mieszka na wysokości, * a na niskość spogląda na niebie i na ziemi?\\
6. Podnosi z ziemi biednego * i z gnoju wywyższa ubogiego:\\
7. Aby go posadzić z książętami, * z książętami ludu swego.\\
8. Który czyni, iż niepłodna mieszka w domu, * matka synów z weselem.

%---------------------
\red{\textbf{Ant.3}} W krzaku, który widział Mojżesz, * płonącym, a nie palącym się, widzimy obraz chwalebnego dziewictwa Twego: Boża Rodzicielko, wstawiaj się za nami.

\red{\textbf{Psalm 121}}\\
1. Uweseliłem się z tego, co mi powiedziano: * Pójdziemy do domu Pańskiego.\\
2. Stały nogi nasze * w dziedzińcach twoich, Jeruzalem.\\
3. Jeruzalem, które się buduje jak miasto, * w sobie zespolone.\\
4. Bo tam wstępowały pokolenia, pokolenia Pańskie, * ustawa dla Izraela, by sławić imię Pana.\\
5. Bo tam stały stolice sądowe, * stolice dla domu Dawidowego.\\
6. Proście o to, co jest ku pokojowi dla Jeruzalem, * a dostatek niech będzie dla tych, którzy cię miłują.\\
7. Niechaj będzie pokój w warowni twojej * i dostatek w basztach twoich.\\
8. Ze względu na braci moich i bliźnich moich * wzywałem pokoju dla ciebie.\\
9. Dla domu Pana, Boga naszego, * szukałem dobra tobie.

%---------------------
\red{\textbf{Ant.4}} Zakwitła różdżka Jessego: * wyszła gwiazda z Jakóba: Dziewica porodziła nam Zbawiciela; chwalimy cię, Boże nasz.

\red{\textbf{Psalm 126}}\\
1. Jeśli Pan nie zbuduje domu, * próżno pracowali, którzy go budują.\\
2. Jeśli Pan nie będzie strzegł miasta, * próżno czuwa ten, który go strzeże.\\
3. Próżno wam przede dniem wstawać, * wstańcie, skoroście usiedli, którzy pożywacie chleba boleści.\\
4. Gdy da umiłowanym swym spanie: * oto dziedzictwem Pańskim synowie, zapłatą owoc żywota.\\
5. Jak strzały w ręku mocarza, * tak synowie utrapionych.\\
6. Błogosławiony człowiek, który napełnił nimi pragnienie swoje: * nie zawstydzi się, kiedy będzie mówił z wrogami swymi w bramie.

%---------------------
\red{\textbf{Ant.5}} Oto Marja porodziła * nam Zbawiciela, którego Jan zoczywszy, zawołał: Oto Baranek Boży, oto, który gładzi grzechy świata, alleluja.

\red{\textbf{Psalm 147}}\\
1. Chwal, Jeruzalem, Pana, * chwal, Syjonie, Boga twego.\\
2. Albowiem umocnił zawory bram twoich, * pobłogosławił synom twym w tobie.\\
3. On użyczył pokoju granicom twoim * i nasyca cię najwyborniejszym zbożem.\\
4. On puszcza słowo swe na ziemię, * bardzo prędko biegnie mowa Jego.\\
5. On daje śnieg jak wełnę, * a mgłę rozsypuje jak popiół.\\
6. Spuszcza lód swój jak bryły, * kto się ostoi przed zimnem Jego?\\
7. Pośle słowo swoje a roztopi je, * wionie wiatr Jego, a pociekną wody.\\
8. On oznajmia słowo swe Jakubowi, * ustawy i przykazania swe Izraelowi.\\
9. Nie uczynił tak żadnemu narodowi, * i nie objawił im przykazań swoich.

%---------------------
\red{\textbf{Capitulum}}\\
Oto Ja posyłam Anioła mojego, a nagotuje drogę przed obliczem mojem. A zarazem przyjdzie do kościoła swego Panujący, którego wy szukacie, i Aniół przymierza, którego wy chcecie.
%---------------------
\red{\textbf{Hymn}}\\
1. Witaj, Gwiazdo morza, * Rodzicielko Boża, * Dziewico nietknięta * Bramo niebios święta.\\
2. Gabrjel zwiastował * Tę cudu godzinę, * Obdarz nas pokojem * Mażąc Ewy winę.\\
3. Zerwij ziemskie pęta * Rzuć światłości blaski, * Niech grzeszni i ślepi * Doznają Twej łaski.\\
4. O, Matko nieś nasze * Prośby i westchnienia * Temu, co zrodzony * Dla świata zbawienia.\\
5. Dziewico przeczysta, * Prosim Cię w pokorze, * Niechaj czyste serca * Mają sługi Boże.\\
6. I bezpieczną drogą * Prowadź nas do celu, * Wtedy szczęście nasze * Znajdziem w Zbawicielu.\\
7. Chwalmy Ojca, Syna * I Ducha Świętego, * Chwalmy wiecznie Boga * W Trójcy Jedynego. * Amen.

%---------------------
\red{\textbf{Versus}}\\
\red{\Vbar.}~Odpowiedź był wziął Symeon od Ducha Świętego.\\
\red{\Rbar.}~Że nie miał oglądać śmierci, ażby pierwéj oglądał Chrystusa Pańskiego.

%---------------------
\red{\textbf{Ant.Magnificat}}~~Starzec * niósł Dziecię, lecz Dziecię starcem rządziło: Oddał cześć Temu, którego Dziewica poczęła, Dziewica zrodziła, Dziewicą zostając po narodzeniu.

\red{\textbf{Magnificat}}\\
1. Wielbij * duszo moja, Pana.\\
2. I rozradował się duch mój: * w Bogu, Zbawicielu moim.\\
3. Iż wejrzał na niskość służebnice swojéj: * albowiem oto odtąd błogosławioną mię zwać będą wszystkie narody.\\
4. Albowiem uczynił mi wielkie rzeczy, który możny jest: * i święte imię jego.\\
5. A miłosierdzie jego od narodu do narodów: * bojącym się jego.\\
6. Uczynił moc ramieniem swojem: * rozproszył pyszne myśli serca ich.\\
7. Złożył mocarze z stolice: * a podwyższył niskie.\\
8. Łaknące napełnił dobrami: * a bogacze z niczem puścił.\\
9. Przyjął Izraela, sługę swego: * wspomniawszy na miłosierdzie swoje.\\
10 Jako mówił do ojców naszych: * Abrahamowi i nasieniu jego na wieki.

%---------------------
\red{\textbf{Oratio}}\\
Wszechmogący wieczny Boże, Majestatu Twego pokornie błagamy, aby, jak Syn Twój Jednorodzony dnia dzisiejszego w postaci naszego ciała do świątyni został przyniesiony, tak abyśmy z sercem oczyszczonem przed Tobą stanąć mogli.
Przez tegoż naszego Pana Jezusa Chrystusa, Twojego Syna: który z Tobą żyje i króluje w jedności Ducha Świętego Bóg, przez wszystkie wieki wieków. Amen.
\footnotetext{Tłumaczenia pochodzą ze strony divinumofficium.com}
\end{document}